\section{结果分析}

\subsection{结果展示}

\begin{figure}[H]
    \centering
    \fbox{\rule{0pt}{4.2cm}\rule{0.82\textwidth}{0pt}}
    \caption{算法流程或结果示意图占位}
    \label{fig:placeholder}
\end{figure}

图~\ref{fig:placeholder} 可在正式论文中替换为算法流程图、网络结构图、误差对比图或敏感性分析图。若结果部分包含多组实验，建议每一组都采用“结果展示 + 原因解释 + 结论提炼”的固定写法，以避免正文层次失衡。

\subsection{结果解释}

在正式论文中，可围绕总成本、准时率、车辆利用率、鲁棒性或预测误差等指标展开分析。对每组结果，建议同时给出数值结论与业务解释，避免出现“有表无文”或“有结论无证据”的情况。
